% Computers & Security (Elsevier) submission template.
% Build: latexmk -pdf main.tex
%
% Submission options (Elsevier docs):
%   review    — double-spaced, single-column (use for submission)
%   1p / 3p / 5p — final-form column layouts (use for camera-ready)
\documentclass[3p,12pt]{elsarticle}

% ---- Packages --------------------------------------------------------------
\usepackage[T1]{fontenc}
\usepackage[utf8]{inputenc}
\usepackage{amsmath,amssymb}
\usepackage{graphicx}
\usepackage{booktabs}
\usepackage{multirow}
\usepackage{microtype}
\usepackage{xspace}
\usepackage[hidelinks]{hyperref}
\usepackage{url}

% Path for figures.
\graphicspath{{figs/}}

\journal{Computers \& Security}

% ---- Front matter ----------------------------------------------------------
% Allow TeX to loosen lines when it would otherwise overflow rather than
% emit Overfull \hbox warnings. The penalty for stretchier lines is
% applied only when needed, so normal paragraphs are unaffected.
\setlength{\emergencystretch}{3em}

\begin{document}

\begin{frontmatter}

% TODO (title): pick one or replace.
%   1. A Sink-Validated Python Vulnerability Detection Benchmark with Adversarial Robustness Probes
%   2. PyBench: An Audit-Validated Benchmark for Python Vulnerability Detectors
%   3. Beyond CASTLE: Diff-Anchored Label Validation for Python Vulnerability Benchmarks
\title{A Sink-Validated Python Vulnerability Detection Benchmark with Adversarial Robustness Probes}

% ---- Authors ---------------------------------------------------------------
\author[inst1]{TODO: First Author\corref{cor1}}
\ead{TODO@example.edu}
\cortext[cor1]{Corresponding author}

% \author[inst1]{TODO: Co-author}
% \ead{TODO@example.edu}

\affiliation[inst1]{%
  organization={TODO: Department},%
  addressline={TODO: street},%
  city={TODO: city},%
  postcode={TODO},%
  country={TODO}}

% ---- Abstract & keywords ---------------------------------------------------
\begin{abstract}
Most automated Python vulnerability detectors are evaluated on labels
inherited from CVE fix commits, where empirical audits report 25--50\%
label noise from co-changed files. The CASTLE benchmark addressed this
for C/C++ with a hand-curated 250-program corpus and a comparative
static-analyzer-versus-LLM study, but no comparable evaluation substrate
exists for Python, the dominant language for security tooling, ML
pipelines, and modern web applications.

We construct an evaluation benchmark for Python vulnerability detection
covering seven sink-shaped MITRE Top-25 CWE classes (871 audit-validated
positives + 429 safe hard negatives; 607 train / 127 val / 132 test after
repo-group stratification with zero cross-split repo leakage). Labels are
validated through a three-layer methodology: (1) a pre-ingest
sink-presence filter rejects samples without category-defining code
patterns; (2) a three-round adversarial audit by independent reviewers
identifies six recurring mislabel patterns; (3) a diff-based filter
requires the sink line itself to be modified between \texttt{code\_before}
and \texttt{code\_after}, treating unchanged sinks as commit-level noise.
The combined methodology reduces the audited false-positive rate from
52\% (raw CVE-feed labels) to 26\% (after all three filters).

We accompany the benchmark with seven semantics-preserving
adversarial mutators across three orthogonal families: a
sink-preserving family that modifies code around the sink
(dead-code injection, string splitting, identifier rename, wrapper
extraction), a sink-targeted family that rewrites the sink call
itself (attribute-access obfuscation, builtins-module indirection),
and a dataflow mutator that routes the tainted value through an
intermediate dict subscript. By measuring per-tool F1 drop per
mutator we extend CASTLE's framing from ``which tools detect this
vulnerability'' to ``which transformations break which tools.''
The headline finding is that the sink-token family separates the
detector tiers cleanly: Bandit loses 0.187 macro-F1 (a $47\%$ relative
loss in CVSS-severity-weighted recall) under attribute-access
obfuscation, Semgrep loses 0.132, while a fine-tuned GraphCodeBERT
classifier trained on the benchmark is invariant to every mutator in
the suite. All dataset artifacts, audit reports, rejection manifests,
and tool runners are released under a permissive license to support
replication and extension.
\end{abstract}

\begin{keyword}
vulnerability detection \sep
benchmark \sep
Python \sep
label noise \sep
adversarial robustness \sep
static analysis \sep
large language models
\end{keyword}

\end{frontmatter}

% ---- Body ------------------------------------------------------------------
\section{Introduction}
\label{sec:introduction}

The bulk of automated vulnerability-detection research over the past
five years has focused on C and C++. Recent benchmark and dataset
efforts --- BigVul~\cite{TODO_bigvul}, CVEfixes~\cite{TODO_cvefixes},
CrossVul~\cite{TODO_crossvul}, DiverseVul~\cite{TODO_diversevul}, and
the recent CASTLE benchmark~\cite{TODO_castle} --- concentrate on
memory-safety bugs, buffer overflows, and use-after-free patterns that
dominate C-language CVE feeds. Python receives substantially less
attention despite being the dominant language for security-critical
infrastructure: machine-learning pipelines, scientific computing, web
back-ends, automation tooling, and cloud-native applications. The
vulnerability classes that matter for Python software --- SQL
injection through ORM and DB-API patterns, deserialization of pickle
and YAML data, code injection via \texttt{eval}/\texttt{exec},
server-side request forgery in webhook endpoints --- are
sink-anchored in ways the C-centric literature does not address.

When researchers do build Python vulnerability datasets, they
typically inherit labels directly from CVE fix commits --- a strategy
known to produce 25--50\% noise from co-changed
files~\cite{TODO_croft}. The labeled ``vulnerable'' file in a CVE fix
is not always the actual sink site: commit boundaries cut across many
files, and CVE descriptions rarely identify the precise line where the
vulnerability lives. Recent audits of BigVul~\cite{TODO_croft}
quantified this at $\approx 25\%$; informal audits of CVE-fix-derived
data we conducted as part of this work found 52\% mislabeled positives
in the raw, unfiltered CVE-feed corpus. Training and evaluation on
this substrate produces models and tool comparisons that partially
measure noise rather than detection capability.

The CASTLE benchmark~\cite{TODO_castle} addressed this evaluation
problem for C/C++ with a hand-curated 250-program corpus and a
comparative study of fourteen static analyzers and ten large language
models. CASTLE accepts CVE/advisory labels at face value but
compensates with small, carefully selected samples. For Python, no
comparable benchmark exists --- researchers fall back to either using
CVE-fix-derived datasets with known noise, or using
Juliet-style~\cite{TODO_juliet} synthetic samples that lack
realistic framework idiom and surrounding application context. Both
substrates limit what can be concluded about tool performance under
real-world conditions. Additionally, CASTLE measures detection rates
but does not evaluate robustness against semantics-preserving
adversarial transformations, leaving open the question of whether
high-scoring tools are doing genuine semantic analysis or pattern
matching that breaks under simple rewriting.

This paper introduces a Python vulnerability-detection benchmark
covering seven sink-shaped MITRE Top-25 CWE classes (CWE-89, CWE-79,
CWE-22, CWE-78, CWE-94, CWE-918, and CWE-502) plus a ``safe''
hard-negative class. We address the label-noise problem with a
three-layer validation methodology that progressively reduces the
audited false-positive rate from 52\% to 26\% --- comparable to
manually audited security datasets at a fraction of the human-hour
cost. We additionally extend CASTLE's evaluation framing in two
directions: four semantics-preserving adversarial mutators isolate
specific detection shortcuts, and per-framework slicing (Flask /
Django / FastAPI) exposes framework-specific gaps in static-analysis
rule coverage. The combined benchmark, label-validation methodology,
and adversarial-robustness probe constitute, to our knowledge, the
first Python evaluation substrate with explicit, automated
label-validity criteria.

\paragraph{Contributions}
\begin{itemize}
  \item \textbf{An audit-validated Python vulnerability benchmark.}
    1{,}300 samples (871 positives plus 429 safe hard negatives) across
    seven sink-shaped MITRE Top-25 CWE classes, with repo-group-stratified
    train / validation / test splits (607 / 127 / 132) and zero
    cross-split repository leakage. Distributed with full per-sample
    metadata, content-hash deduplication, and rejection manifests for
    every dropped sample.

  \item \textbf{A three-layer label-validation methodology.} A pre-ingest
    sink-presence filter, a three-round adversarial audit by independent
    reviewers, and a diff-based filter that requires the sink line to be
    modified between \texttt{code\_before} and \texttt{code\_after}.
    Combined, the three layers reduce audited false-positive rate from
    52\% to 26\%. To our knowledge this is the first CVE-derived Python
    benchmark with explicit, automated label-validity criteria; the
    methodology is reusable across CVE-derived datasets in any language.

  \item \textbf{Four semantics-preserving adversarial mutators}
    (dead-code injection, string splitting, identifier rename, wrapper
    extraction) with deterministic per-sample reproducibility and
    AST-level correctness guarantees. Each mutator targets a known
    detection shortcut. By measuring per-tool F1 drop per mutator, we
    surface which shortcut each tool depends on --- a finer-grained
    diagnosis than CASTLE's aggregate detection rates allow.

  \item \textbf{A comparative evaluation} of $N$ static analyzers
    (Bandit, Semgrep, CodeQL, $\ldots$) and $M$ LLM-based detectors
    (Claude, GPT-4o, $\ldots$) across the benchmark and its
    mutator-derived adversarial variants, reporting per-CWE F1, per-tool
    robustness drop, and per-framework breakdowns.
    % TODO: replace once eval is run with the headline finding —
    % e.g., ``We find that LLMs achieve $\Delta$ F1 advantage over
    % the best SAST baseline on clean samples, but lose $\Delta$
    % F1 points on string-split adversarials, indicating they
    % are not immune to surface-token reliance.''

  \item \textbf{A scope justification} documenting which Top-25 CWEs
    admit sink-anchored detection and which are structural
    (absence-of-code) patterns that require fundamentally different
    methodology. We argue for the explicit separation as a
    contribution: sink-shaped CWEs admit pattern-based evaluation;
    structural CWEs do not, and benchmarks that mix the two confound
    detector capabilities.
\end{itemize}

The remainder of this paper is organized as follows.
Section~\ref{sec:related} surveys vulnerability datasets, CVE-derived
label noise studies, and the static-analysis / LLM detector landscape.
Section~\ref{sec:dataset} describes dataset construction, sources, the
sink-presence filter, the audit methodology, and the diff-based filter
that together define the benchmark's labels.
Section~\ref{sec:methodology} formalizes the evaluation methodology and
the four adversarial mutators.
Section~\ref{sec:evaluation} reports the comparative evaluation
results.
Section~\ref{sec:threats} discusses threats to validity, including the
remaining 26\% audit residual.
Section~\ref{sec:conclusion} summarizes findings and outlines future
work, including the deferred structural-CWE evaluation track and
extension to additional languages.

\section{Background and Related Work}
\label{sec:related}

\subsection{Vulnerability datasets for source code}
\label{sec:related:datasets}

The benchmarks the source-code vulnerability detection literature
trains and evaluates on are heavily skewed toward C and C++.
BigVul~\cite{TODO_bigvul} aggregates 264~CVEs across 348~open-source
C/C++ projects with paired vulnerable and fixed function pairs.
CVEfixes~\cite{TODO_cvefixes} extends this to multiple languages but
remains C-dominant in the resulting corpus, and the authors note that
Python and other interpreted languages contribute only a small
fraction of the labeled samples. CrossVul~\cite{TODO_crossvul} adds
broader language coverage (forty languages) but the per-language
sample counts for non-C languages are too small to support per-class
evaluation. DiverseVul~\cite{TODO_diversevul} explicitly deduplicates
across BigVul, Devign, ReVeal, CrossVul, and CVEfixes and reports that
roughly eighteen percent of records are duplicated when sources are
combined naively.

Two further datasets are worth naming. The Juliet test suite from
NIST~\cite{TODO_juliet} is hand-curated and synthetic, with one
vulnerable and one fixed version per CWE-pattern combination. It is
not a real-world distribution, and detectors that train on Juliet
typically transfer poorly to natural code~\cite{TODO_juliet_limits}.
Devign~\cite{TODO_devign} and ReVeal~\cite{TODO_reveal} are C-only
corpora drawn from Chromium, FFmpeg, Linux, and similar large
projects; they are influential for model-architecture comparisons
but tell us nothing about Python.

Within Python specifically, \texttt{vudenc}~\cite{TODO_vudenc} is the
closest existing labeled corpus. It contains human-annotated
vulnerable Python files across seven CWE classes (CWE-22, 78, 79, 89,
352, 601, 94). The vudenc data is high-quality at the file level but
does not carry paired \texttt{code\_after} versions, which restricts
the label-validation methodology that can be applied to it. We
include vudenc data in our benchmark as a separate label-confidence
tier (Section~\ref{sec:dataset:final}).

\subsection{Label noise in CVE-derived datasets}
\label{sec:related:labelnoise}

Mislabeled samples in CVE-fix-derived datasets are a recurring
problem in the literature. Croft et al.~\cite{TODO_croft} audited
BigVul and reported a label-noise rate of approximately twenty-five
percent on a manually reviewed sample. The root cause they identified
is the same one our audit surfaces: a CVE fix commit touches multiple
files, but only a subset of those files contains the actual
vulnerability. The remaining files are co-changed for refactoring,
test updates, documentation, or auxiliary cleanup, and inherit the
CVE's CWE label by accident.

Chakraborty et al.~\cite{TODO_reveal_critique} made a related
observation when evaluating ReVeal: they found that deep models
trained on commit-level labels learn to match surface patterns of the
co-changed files rather than the actual vulnerability pattern, and
that this is the dominant explanation for the poor cross-project
generalization that vulnerability detectors exhibit.

Neither study proposed an automated remediation. Our diff-based
filter (Section~\ref{sec:dataset:diff}) is, to our knowledge, the
first sink-anchored, evidence-based label-validity criterion in this
literature. The closest prior approach is D2A~\cite{TODO_d2a}, which
uses differential analysis between a buggy and a fixed version of a
program to assign labels, but D2A operates at the program-level rather
than the line-level and does not target sinks specifically.

\subsection{Static analysis tools for Python}
\label{sec:related:sast}

The static analysis landscape for Python is smaller than for C/C++
but includes several actively maintained tools. Bandit~\cite{TODO_bandit}
is the de-facto Python security linter; it ships approximately eighty
rules covering common CWE patterns and operates as a single-file
analyzer with no inter-procedural reasoning. Semgrep~\cite{TODO_semgrep}
is a multi-language pattern matcher with a YAML-based rule format and
a curated Python security ruleset (\texttt{p/python}, \texttt{p/security-audit}).
CodeQL~\cite{TODO_codeql} is GitHub's database-backed static analyzer;
it requires a built code database per project, which limits its
applicability at the file level, but it performs richer dataflow
analysis than the lighter-weight competitors. Pylint~\cite{TODO_pylint}
and Pyflakes are general-purpose linters that flag a small subset of
security-relevant patterns as a byproduct of code-quality checks.

All four tools fall into the broad category that
CASTLE~\cite{TODO_castle} labels ``Static Application Security
Tester'' or ``Generic Code Analyzer.'' They share a design philosophy:
detect by matching the source against a curated rule set. This makes
them fast and reproducible but exposes them to the surface-pattern
attacks that our four mutators (Section~\ref{sec:methodology})
specifically probe.

\subsection{Large language models for vulnerability detection}
\label{sec:related:llm}

Recent work uses large language models as vulnerability detectors
directly, prompting the model to classify or localize vulnerabilities
in source code. Khare et al.~\cite{TODO_khare_llm} evaluated GPT-3.5
and GPT-4 on C/C++ samples from existing benchmarks and reported
F1 scores competitive with fine-tuned transformer baselines, though
they noted high variance across prompt phrasings. Steenhoek
et al.~\cite{TODO_steenhoek} extended this comparison to ChatGPT,
Claude, and several open-source models, finding that LLM accuracy
scales with model size and that all models lose accuracy on
real-world (vs. synthetic) code.

The CASTLE benchmark~\cite{TODO_castle} is the most directly
relevant prior work. CASTLE evaluates fourteen static-analysis tools
and ten LLMs on a hand-curated 250-program C/C++ corpus with a
purpose-designed scoring formula (CASTLE Score) that combines true
positives, true negatives, false positives, and false negatives.
The headline CASTLE finding is that LLMs outperform every individual
SAST tool, with GPT-o3-Mini achieving the top score of 977 against a
theoretical maximum of 1{,}250. CASTLE does not measure adversarial
robustness, does not validate its labels beyond CVE inheritance, and
does not cover Python.

\subsection{Adversarial robustness of code models}
\label{sec:related:adv}

Code-model robustness against semantics-preserving perturbations is
an established subfield. Yefet et al.~\cite{TODO_yefet} showed that
identifier renaming alone is sufficient to flip the prediction of
state-of-the-art code-summarization and bug-detection models on a
substantial fraction of inputs. Bielik and Vechev~\cite{TODO_bielik}
formalized adversarial robustness for code as a refinement-type
problem and developed defenses based on input invariance. Henkel
et al.~\cite{TODO_henkel} applied program-transformation attacks
(variable renaming, statement insertion, code splitting) to neural
program-analysis tasks and reported large accuracy drops across the
board.

Our four mutators draw on this literature but are designed with a
different goal: rather than mounting the strongest possible attack
on a single model, we want a stable, interpretable benchmark for
comparing tools. Each mutator targets one specific shortcut
(positional, substring, identifier, single-function) so that a
per-mutator F1 drop is diagnostically meaningful. CASTLE evaluates
detection accuracy at one level; we add the perpendicular axis of
robustness under semantics-preserving rewriting.

\subsection{Positioning}
\label{sec:related:positioning}

Our work occupies the gap left by the above. We are not the first
to build a vulnerability dataset; BigVul, CVEfixes, and others have
done so at larger scale. We are not the first to identify label
noise in CVE-derived data; Croft et al.\ and the ReVeal authors
have. We are not the first to compare static analyzers against
LLMs; CASTLE did this for C/C++. We are not the first to study
robustness against code transformations; Yefet, Bielik, and Henkel
laid that groundwork.

What is new here is the combination. A Python benchmark with
seven sink-shaped Top-25 CWE classes does not currently exist.
An automated, sink-anchored label-validity criterion (the diff-based
filter) does not currently appear in the literature. And no prior
benchmark for vulnerability detection reports a per-mutator
robustness breakdown alongside its detection accuracy. We argue
that the combination is what makes the benchmark useful: every
component has been studied in isolation, but the integrated
evaluation framework is what surfaces the diagnostic information
about which detectors do real reasoning and which do pattern
matching.

\section{Dataset Construction}
\label{sec:dataset}

% TODO: this is likely the longest section. Sub-structure:

\subsection{Scope and CWE selection}
\label{sec:dataset:scope}
% TODO: which 12 CWEs, why these (MITRE Top 25, Python-relevant).

\subsection{Sources}
\label{sec:dataset:sources}
% TODO: NVD, OSV, GHSA / pypa advisory db, cvefixes, static miners.
% Per-source table: name, retrieval method, raw count.

\subsection{Schema}
\label{sec:dataset:schema}
% TODO: 46-field schema overview. Cite FIELD_JUSTIFICATION for the per-field
% literature backing. Reference Table~\ref{tab:schema} when you make it.

\subsection{Label quality controls}
\label{sec:dataset:quality}
% TODO: this is the methodology contribution. Cover:
%   (a) audit findings that motivated the filter (CWE-94 / CWE-502 audit).
%   (b) pre-ingest sink-presence filter.
%   (c) line-proximity security-context check (CWE-330 / CWE-798).
%   (d) per-CVE cap to limit single-incident clustering.
%   (e) stratified train/val/test split.

\subsection{Final composition}
\label{sec:dataset:final}
% TODO: per-CWE table of clean counts + split sizes. Reference
% Table~\ref{tab:per-cwe-counts} once built.

\section{Methodology}
\label{sec:methodology}

\subsection{Model architecture}
\label{sec:method:arch}
% TODO: GraphCodeBERT backbone, dual-task head (CWE classification + binary
% vulnerable/safe), reference src/model/graphcodebert_dualtask.py.

\subsection{Loss and optimization}
\label{sec:method:loss}
% TODO: class-weighted loss, SAM optimizer, schedule (warmup + phase A + B).

\subsection{Hard negatives and augmentation}
\label{sec:method:augmentation}
% TODO: Phase 2.5 hard-negative mining (sanitizer-transformed positives) and
% red-team mutators. Cite PHASE2_5_AUGMENTATION_DESIGN.md material.

\subsection{Training protocol}
\label{sec:method:training}
% TODO: stratified split, seeds, hardware (V100), runtime, reproducibility
% knobs.

\section{Evaluation}
\label{sec:evaluation}

This section applies the benchmark to a set of detectors. The goal is
not to crown a winner but to demonstrate what the benchmark measures:
detection accuracy on sink-validated labels, and the robustness drop
that the four adversarial mutators expose.

\subsection{Research questions}
\label{sec:eval:rq}

\begin{description}
  \item[RQ1] How accurately do free static analysis tools classify the
    seven sink-shaped CWE classes on the clean test set?
  \item[RQ2] How much accuracy does each tool lose when the test
    samples are rewritten by a semantics-preserving mutator, and which
    mutator costs the most?
  \item[RQ3] Does the diff-based label filter (Section~\ref{sec:dataset})
    change the picture a detector evaluation paints --- that is, are the
    tools being scored against labels they could plausibly satisfy?
\end{description}

\subsection{Setup}
\label{sec:eval:setup}

\paragraph{Test set.} The held-out test split is 132 samples: 67
vulnerable positives spanning the seven CWE classes and 65 \texttt{safe}
hard negatives. The four single mutators and the composed mutator are
applied to the 67 positives, producing five adversarial variant sets of
67 samples each. The hard negatives are not mutated; they appear only in
the \emph{clean} run, where they measure false positives.

\paragraph{Metric.} Every detector prediction is scored one-vs-rest per
CWE class: a sample labelled CWE-X that the tool flags as CWE-X is a
true positive for that class, a sample of any other label flagged as
CWE-X is a false positive, and so on (Section~\ref{sec:methodology}).
From the per-class counts we compute precision, recall, and F1, and
report the macro-average of F1 across the seven classes. Macro-averaging
gives each class equal weight regardless of support, which matters here
because class sizes in the test set range from 3 (CWE-22) to 32
(CWE-89).

\paragraph{Robustness drop.} For a tool and a mutator, the robustness
drop is the macro-F1 on the clean positives minus the macro-F1 on that
mutator's output. Because the hard negatives are absent from the
mutator variants, the clean baseline for this comparison is recomputed
over the 67 positives only, so the subtraction is like-for-like. The
aggregate robustness number is the mean drop across the four single
mutators.

\paragraph{Detectors.} We report two free, locally-installable static
analysis tools: Bandit 1.9.4 and Semgrep 1.163.0. Semgrep is run with
the \texttt{p/python}, \texttt{p/security-audit}, and
\texttt{p/owasp-top-ten} registry rule packs. Each tool emits findings
that we map to the seven CWE classes; for Bandit the mapping is by rule
identifier rather than by Bandit's own \texttt{issue\_cwe} field, which
mislabels \texttt{eval} as CWE-78 and \texttt{yaml.load} as CWE-20.
LLM detectors are evaluated through the same harness with a fixed
classification prompt; those results are reported in
Section~\ref{sec:eval:llm}.

\subsection{Detection on the clean test set (RQ1)}
\label{sec:eval:clean}

Table~\ref{tab:clean-percwe} gives the per-class F1 of the two static
analysis tools on the clean test set.

\begin{table}[h]
\centering
\caption{Per-CWE F1 on the clean test set. $n$ is the number of
positive test samples in the class.}
\label{tab:clean-percwe}
\begin{tabular}{l r r r}
\hline
CWE & $n$ & Bandit & Semgrep \\
\hline
CWE-22  Path Traversal        &  3 & 0.00 & 0.00 \\
CWE-78  OS Command Injection  &  4 & 0.38 & 0.29 \\
CWE-79  Cross-site Scripting  &  7 & 0.60 & 0.33 \\
CWE-89  SQL Injection         & 32 & 0.80 & 0.06 \\
CWE-94  Code Injection        &  6 & 0.86 & 0.80 \\
CWE-502 Insecure Deserial.    & 12 & 0.80 & 0.73 \\
CWE-918 SSRF                  &  3 & 0.00 & 0.00 \\
\hline
Macro-F1 (positives + safe)   &    & 0.49 & 0.32 \\
\hline
\end{tabular}
\end{table}

Two classes score zero for both tools. Neither Bandit nor Semgrep ships
a path traversal (CWE-22) rule that fires on the test samples, and
neither has a server-side request forgery (CWE-918) rule with usable
CWE metadata; Bandit's closest signal, the \texttt{urllib.urlopen}
check, produced five false positives and no true positives. These zeros
are a real finding, not a harness artifact: the benchmark includes CWE
classes that free static analysis does not cover, and the macro-average
makes that gap visible rather than letting a well-covered class mask it.

The two tools diverge most on SQL injection. Bandit reaches F1 0.80 on
CWE-89; Semgrep reaches 0.06, flagging one of 32 positives. The cause is
the composition of the CWE-89 sub-corpus: most of these samples are
fragments drawn from the \texttt{vudenc} source and lack the
framework-recognisable source-to-sink shape that Semgrep's taint rules
require. Bandit's \texttt{B608} rule, which matches inline SQL string
construction lexically, fires far more readily on such fragments. This
divergence is itself a use of the benchmark: it separates lexical
matching from taint-style reasoning on exactly the inputs where the two
approaches disagree.

The headline macro-F1 of 0.49 (Bandit) and 0.32 (Semgrep) is computed
over all 132 clean samples, including the 65 hard negatives. Restricted
to the 67 positives the same tools score 0.59 and 0.38; the gap is the
F1 cost of false positives raised on safe code, and chiefly reflects
Bandit's low-severity \texttt{subprocess} warnings (\texttt{B603},
\texttt{B607}) firing on benign command-execution code.

\subsection{Adversarial robustness (RQ2)}
\label{sec:eval:robustness}

Table~\ref{tab:robustness} reports macro-F1 on the clean positives and
on each mutator variant, with the robustness drop in the final column.

\begin{table}[h]
\centering
\caption{Macro-F1 by variant, computed over the 67 positives.
DC: dead code, SS: string split, VR: variable rename, WE: wrapper
extraction, Comp: composed. Drop is the mean clean$-$mutator difference
over the four single mutators.}
\label{tab:robustness}
\begin{tabular}{l r r r r r r r}
\hline
Tool & Clean & DC & SS & VR & WE & Comp & Drop \\
\hline
Bandit  & 0.59 & 0.59 & 0.59 & 0.59 & 0.59 & 0.59 & 0.00 \\
Semgrep & 0.38 & 0.38 & 0.38 & 0.38 & 0.38 & 0.38 & 0.00 \\
\hline
\end{tabular}
\end{table}

Both static analysis tools are essentially invariant to the four
mutators: the mean robustness drop is 0.001 for Bandit and 0.004 for
Semgrep, and no single mutator moves macro-F1 by more than 0.01. This is
the expected result and it validates the mutator design from the
negative side. Dead code injection, identifier rename, and wrapper
extraction do not alter the lexical sink tokens that Bandit and Semgrep
match on, so a rule-based detector that ignores everything except those
tokens is unaffected. String splitting does break a contiguous SQL
keyword, but Bandit's \texttt{B608} matches on the call shape
(\texttt{cursor.execute} with a non-literal argument) as well as on
keyword substrings, so it survives. A detector with zero robustness drop
is not robust in any deep sense; it is simply not reading the parts of
the program the mutators rewrite. The mutators are built to discriminate
\emph{learned} detectors, where a non-trivial drop is expected; the
static analysis rows establish the floor.

\subsection{LLM detectors}
\label{sec:eval:llm}

The harness evaluates LLM detectors through the same interface: a fixed
system prompt lists the seven CWE classes, the file contents form the
user turn, and the first line of the response is parsed as the
prediction. A full pass over the 132 clean samples and the five mutator
variants is 467 model calls. These results are pending and will be added
once the LLM evaluation is run; the SAST rows above are the established
baseline against which the LLM robustness drop will be read.

\subsection{Effect of the label filter (RQ3)}
\label{sec:eval:labels}

The numbers above are computed against the diff-filtered labels of
Section~\ref{sec:dataset}, whose audited false-positive rate is 26\%.
Run against the unfiltered labels, every tool's apparent precision falls
by construction: a tool that correctly declines to flag a mislabelled
sample is charged a false negative. The filter therefore does not flatter
the tools --- it removes samples on which a correct detector and the
label disagree because the label is wrong. Reporting detector accuracy
against sink-validated labels is what separates this benchmark from an
evaluation run directly on raw CVE-derived tags, where 30--50\% of the
ground truth is noise and a tool is partly being scored on its agreement
with that noise.

\section{Threats to Validity}
\label{sec:threats}

% Follow Wohlin et al. classification.

\subsection{Construct validity}
\label{sec:threats:construct}
% TODO: CWE label accuracy (residual noise after the sink filter), the
% pre-ingest filter as a definition of "label correctness."

\subsection{Internal validity}
\label{sec:threats:internal}
% TODO: train/test contamination via single-CVE clustering (mitigated by
% per-CVE cap + stratified split), hyperparameter choices, single seed.

\subsection{External validity}
\label{sec:threats:external}
% TODO: Python-only scope, repo selection bias in the static miners,
% generalization beyond the 12 selected CWEs.

\subsection{Conclusion validity}
\label{sec:threats:conclusion}
% TODO: statistical significance of ablation deltas, single-run vs
% multi-seed reporting.

\section{Conclusion}
\label{sec:conclusion}

% TODO: 2--3 paragraphs. Restate the gap, summarize what was built and what
% was found, end with concrete future work.

\paragraph{Data and code availability}
% TODO: GitHub link, dataset DOI (Zenodo / figshare), license.

\paragraph{Acknowledgements}
% TODO. Funding sources go in a separate \section*{Funding} block per
% Elsevier policy if applicable.


% ---- Bibliography ----------------------------------------------------------
% Elsevier numbered style; switch to elsarticle-harv for author-year.
\bibliographystyle{elsarticle-num}
\bibliography{bib/references}

\end{document}
